\subsection{The Cauchy criterion for a function limit}

Heine's theorem translates a function limit into convergence of output sequences. The corresponding Cauchy criterion removes the need to name the candidate limit in advance.

\begin{theorembox}[Cauchy criterion for a finite function limit]
Let $f\colon D\to\mathbb R$, where $D\subseteq\mathbb R$, and let $a$ be an accumulation point of $D$. A finite limit
\[
\lim_{x\to a}f(x)
\]
exists if and only if, for every $\varepsilon>0$, there is $\delta>0$ such that whenever $x,y\in D$ satisfy
\[
0<|x-a|<\delta,
\qquad
0<|y-a|<\delta,
\]
we have
\[
|f(x)-f(y)|<\varepsilon.
\]
\end{theorembox}

\begin{center}
\begin{tikzpicture}[x=1.15cm,y=0.95cm,>=Latex]
  \def\aa{0.4}
  \def\del{1.35}
  \def\xx{-0.65}
  \def\yy{1.25}
  \def\fx{1.58}
  \def\fy{2.34}

  \fill[blue!9] ({\aa-\del},-0.35) rectangle ({\aa+\del},3.8);
  \draw[->] (-2.4,0) -- (3.35,0) node[right] {$t$};
  \draw[->] (0,-0.35) -- (0,3.8) node[above] {$f(t)$};

  \draw[dashed,blue!70!black] ({\aa-\del},-0.2) -- ({\aa-\del},3.65);
  \draw[dashed,blue!70!black] ({\aa+\del},-0.2) -- ({\aa+\del},3.65);
  \draw[thick,domain=-2.1:3.0,samples=80] plot (\x,{0.4*\x+1.84});

  \draw[dashed] (\xx,0) -- (\xx,\fx);
  \draw[dashed] (\yy,0) -- (\yy,\fy);
  \draw[dashed] (0,\fx) -- (2.75,\fx);
  \draw[dashed] (0,\fy) -- (2.75,\fy);

  \fill[blue!75!black] (\xx,\fx) circle (2.5pt);
  \fill[blue!75!black] (\yy,\fy) circle (2.5pt);
  \node[above left] at (\xx,\fx) {$(x,f(x))$};
  \node[above right] at (\yy,\fy) {$(y,f(y))$};

  \draw[<->,thick,orange!85!black] (2.7,\fx) -- (2.7,\fy)
    node[midway,right=4pt] {$|f(x)-f(y)|<\varepsilon$};
  \draw[<->,thick,blue!70!black]
    ({\aa-\del},-0.27) -- ({\aa+\del},-0.27)
    node[midway,below=3pt] {$x,y\in(a-\delta,a+\delta)$};

  \node[below] at (\aa,0) {$a$};
  \node[below] at (\xx,0) {$x$};
  \node[below] at (\yy,0) {$y$};
\end{tikzpicture}
\end{center}

\begin{interpretationbox}[The Cauchy picture]
The criterion does not begin by naming a target value $L$. It asks whether any two outputs obtained from sufficiently near admissible inputs must be close to each other. Shrinking the input neighbourhood must make the entire collection of nearby outputs fit into an arbitrarily small vertical spread.
\end{interpretationbox}

\begin{remarkbox}[Proof status]
The criterion is the function analogue of the Cauchy criterion for sequences and ultimately uses completeness of $\mathbb R$. We use the equivalence and explain its meaning, but do not repeat the full completeness argument.
\end{remarkbox}

The criterion says that all function values sufficiently near $a$ must be close to one another. It does not initially specify the number to which they converge.

\begin{examplebox}[A jump fails the Cauchy test]
Let
\[
f(x)=
\begin{cases}
-1,&x<0,\\
1,&x>0.
\end{cases}
\]
No matter how small $\delta$ is, we can choose $x<0<y$ with
\[
|x|<\delta,
\qquad
|y|<\delta.
\]
Then
\[
|f(x)-f(y)|=|-1-1|=2.
\]
The nearby values do not become mutually close, so the limit at $0$ does not exist.
\end{examplebox}

\begin{interpretationbox}[Heine and Cauchy describe the same settling]
Heine follows every approaching input sequence. Cauchy compares nearby output values directly. Both formulations express that all admissible approaches must settle into one common behaviour.
\end{interpretationbox}
